% !TEX root = "./main.tex"
% Meta-Daten, die mehrfach verwendet werden können
% Angaben Autor*in
\newcommand{\AuthorFirstName}{Max}
\newcommand{\AuthorLastName}{Mustermann}
\newcommand{\AuthorStudentNumber}{1234567}
\newcommand{\AuthorStreet}{Musterstraße 1}
\newcommand{\AuthorZipCode}{12345}
\newcommand{\AuthorLocality}{Musterstadt}
\newcommand{\AuthorEmail}{max.mustermann@smail.inf.h-brs.de}
% Kombinationen
\newcommand{\Author}{\AuthorFirstName\ \AuthorLastName}
\newcommand{\AuthorAddress}{\AuthorStreet \newline \AuthorZipCode\ \AuthorLocality}
% Für externe Arbeiten
\newcommand{\ExternalCompany}{Musterfirma}
\newcommand{\Location}{Musterstadt}
% Angaben zum Studium
\newcommand{\CourseOfStudies}{Musterstudium}
\newcommand{\CourseOfStudiesDegree}{Bachelor / Master}
% Angaben zum Dokument
\newcommand{\DocumentType}{Thesis}
\newcommand{\DocumentTitle}{Mustertitel}
\newcommand{\DocumentSubject}{\DocumentType im \CourseOfStudiesDegree\-Studiengang \CourseOfStudies}
\newcommand{\Keywords}{}
\newcommand{\FirstSupervisor}{Prof. Dr. Martina Mustermann}
\newcommand{\SecondSupervisor}{Prof. Dr. Maria Mustermann}
% Entweder festes Datum oder \today
\newcommand{\DateOfSubmission}{\today}
% Eigentlich obsolet, aber vielleicht kann man das ja auch für andere Fachbereiche benutzen.
\newcommand{\Department}{Computer Science}
\newcommand{\DepartmentGER}{Informatik}